\section{Threats to validity}
\label{sec:threats}

We follow the usual classification of threats~\cite{wohlin2012}.

\paragraph{Internal validity} The authors designed the artefact and evaluated it. Three measures limit the bias this invites: every protocol, with its parameters and expected outcomes, was committed to the public repository before the measurement it governs; every acceptance criterion is checked by an automated assertion; and every result file is public. Two deviations are reported as such. The first cost run and the first anchoring run at 2\,s ran beside unrelated containers. Both were repeated with nothing else running, and the first files are kept and marked as superseded. A review of the manuscript found that the redacted export kept a policy digest from which the thresholds could be recovered. The exporter was changed and the audit questions were run again under an amendment recorded beforehand. The first of these runs matched a metric value inside a timestamp and was repeated, and its files are also kept as superseded. The analysis of stability across 500 partitions was decided after the main run and is reported as post hoc.

\paragraph{External validity} The evaluation covers one regulated domain, credit scoring, and one industrial stream, and the drift in the credit case was induced. The Fabric network is a laboratory network: two organisations on one machine, one Raft orderer, no TLS between chaincode and peer, and demonstration keys. All timings come from one machine, an Apple M4 Pro, on which \path{fsync} does not force the drive's write cache. This favours neither arm of the ablation, but absolute latencies on a server with durable writes would be higher. The MultiFlow streams were replayed from the platform's files without Kafka, although the replay reproduced the live runs batch by batch.

\paragraph{Construct validity} The matrix of Table~\ref{tab:matrix} rests on the authors' reading of the obligations and has not been reviewed by an external legal expert. The fourteen audit questions derive from the legal text, but the authors chose and operationalised them, and in configuration C a signed log anchored after every record stands in for the ledger in all questions but Q12. No auditor, notified body or authority evaluated the artefact, so its usefulness to them is not claimed.

\paragraph{Conclusion validity} Anchoring was measured with serial submission, one transaction in flight and three repetitions per setting. Concurrent clients would change throughput and latency. The comparison of \path{VerifyChain} across state databases rests on chains of different length and is reported per record only. The paired differences of the ablation are tens of microseconds and their interquartile ranges include zero for the signature, so the article reports medians of round medians over 2,000 paired observations and does not claim a difference smaller than these allow.
